% Simplified Chinese labels and lesson macros.
\renewcommand{\contentsname}{目录}
\renewcommand{\figurename}{图}
\renewcommand{\tablename}{表}
\hypersetup{pdftitle={从零理解 simpleKafka},pdfauthor={simpleKafka 教学项目}}

\newenvironment{contract}{\begin{quote}\small\noindent\textbf{行为契约与不变量。}\ }{\end{quote}}
\newcommand{\lessoncommand}[1]{\par\noindent\textbf{命令与晋级。} \coursecmd{./gradlew :stepTest -Pstep=#1} 用于定位本步；\coursecmd{./gradlew :test -Pstep=#1} 从第 1 步累计运行至本步。仅当本章末尾编号（4、8、12、16、20、24、28、31）的累计测试全绿，才算完成本章并进入下一章。}
\newcommand{\lessonquestions}[2]{\subsubsection*{自检问题}\begin{enumerate}\item #1\item #2\end{enumerate}}
\newcommand{\lessonhints}[3]{\subsubsection*{分级提示（不是答案）}\begin{hintlevels}\item[一级] #1\item[二级] #2\item[三级] #3\end{hintlevels}}
\newcommand{\lessonsection}[2]{\section{第 #1 步：#2}\label{step:#1}}
\newcommand{\stepref}[2]{\hyperref[step:#1]{第 #1 步（#2）}}
\newcommand{\expected}[1]{\begin{quote}\small\textbf{可观测例子与测试点。} #1\end{quote}}
\newcommand{\providededit}[2]{\begin{quote}\small\textbf{起点与编辑位置。}\par 已给：#1\par 请修改：#2\end{quote}}
\newcommand{\pitfalls}[1]{\subsubsection*{常见误区与真实 Kafka 的差异}#1}
\newcommand{\goals}[2]{\subsubsection*{目标与前置知识}本步目标：#1\par\textbf{开始前：}#2}
\newcommand{\background}[1]{\subsubsection*{背景与推理}#1}
\newcommand{\lessoncontract}[1]{\begin{contract}#1\end{contract}}
